\section{Axiom Inheritance and the Non-Self-Balancing Finding}
\label{sec:axiom_inheritance}

Since $\volRlower$ is a lower bound on $\volR$, and $\volR =
\volL$ restricted to the exercised sub-poset, the axiom
inheritance analysis applies directly, with an additional layer:
the lower bound itself is not a measure (it is a bound on a
measure).

\begin{theorem}[Axiom-Analogue Analysis for $\volRlower$]
\label{thm:axiom_inheritance}
$\volRlower$ is a lower-bound statistic, not a poset measure in
the sense of~\citep[Proposition~1]{lasser2026poset}; the clauses
below describe \emph{operational analogues} of M1--M6 for
$\volRlower$, not literal axiom inheritance.  M4 and the
structural M6-failure result are substantive; M1, M2, M3, and
M5 are weaker conditional statements noted for completeness.

\begin{description}
\item[(M1) Non-negativity (analogue).] Satisfied.  $\volRlower$ is
  a sum of non-negative terms ($\Delta\volL(X_n) \geq 0$ per event;
  max-attribution preserves non-negativity).

\item[(M2) Null empty set (analogue).] Satisfied.  With zero
  sacrifice events, $\volRlower = 0$.

\item[(M3) Monotonicity (conditional analogue).] Satisfied in
  the observational direction: adding admissible sacrifice
  events weakly increases the constructed lower bound under
  Part~A's sub-batch-supremum form
  ($\volRlower^{\mathrm{A}}$,
  Proposition~\ref{prop:monotone_part_a}) and under Part~B's
  per-capability max-attribution form
  ($\volRlower$,
  Proposition~\ref{prop:monotone}).  Both forms inherit
  monotonicity, with Part~A applied to the sub-batch supremum
  over admissible event subsets and Part~B applied to the
  per-capability max over events covering each capability.
  M3 does \emph{not} hold
  for the relationship between $\volL$ and $\volRlower$: adding
  a benchmarked capability to the poset increases $\volL$ but
  does not affect $\volRlower$ unless that capability is
  subsequently sacrificed.  This is a lower-bound statistic
  property, not a poset-measure axiom.

\item[(M4) Additivity under disjointness (substantive).]
  Satisfied under the parent-poset hypothesis.  By construction,
  $\volR$ is $\volL$ restricted to the exercised sub-poset
  $P^{\mathrm{ex}} \subseteq P$.  M4 for $\volL$ on $P$ states
  that $\volL$ is additive over components of $P$ that are
  pairwise poset-disjoint (no subsumption or
  cooperative-composition links).  Take subsets $A, B
  \subseteq P^{\mathrm{ex}}$ that are
  \emph{poset-disjoint in $P$}---neither cooperative nor
  subsumption edges in the full poset connect any $a \in A$ to
  any $b \in B$, even via unexercised mediating capabilities.
  Note this is a stronger hypothesis than poset-disjointness in
  $P^{\mathrm{ex}}$ alone: the latter could permit $a, b$ to
  share a parent-poset cooperative whose other participants
  happen to be unexercised in $W$, in which case $A, B$ would
  be poset-connected in $P$ even though neither cooperative
  nor subsumption edges restricted to $P^{\mathrm{ex}}$
  connect them.  Under poset-disjointness in $P$, M4 on $P$
  gives $\volL(A \cup B) = \volL(A) + \volL(B)$; restricting
  both sides to $P^{\mathrm{ex}}$ gives $\volR(A \cup B) =
  \volR(A) + \volR(B)$, establishing M4 for $\volR$.
  Theorem~\ref{thm:aggregate_bound}'s pairwise
  poset-independence hypothesis is stated against the parent
  poset $P$ for this reason; the aggregate bound's M4
  application uses parent-poset independence, not
  $P^{\mathrm{ex}}$-internal independence.

\item[(M5) Non-triviality (conditional analogue).] A single
  sacrifice event with $\Delta\volL(X) > 0$ produces
  $\volRlower > 0$ on the capabilities the event covers.  This
  is an observation-dependent analogue, not the structural
  non-triviality of every valid capability.

\item[(M6) Superadditivity under independence.]
  The situation splits into a measure-axiom layer and a
  system-level layer:
  \begin{enumerate}
  \item[(a)] \textbf{$\volR$ as a measure on a fixed exercised
    sub-poset inherits M6.}
    For any fixed $P^{\mathrm{ex}}_W$, if $A, B \subseteq P$ are
    poset-independent in $P$, then $A \cap P^{\mathrm{ex}}_W$
    and $B \cap P^{\mathrm{ex}}_W$ are poset-independent in
    $P^{\mathrm{ex}}_W$, and $\volL$'s M6 on $P$ lifts to
    $\volR$'s M6 on $P^{\mathrm{ex}}_W$ by the same restriction
    argument used for M4 above.  This is literal axiom
    inheritance for the set function $\volR$.

    The structurally relevant failure is at the
    \emph{system level}, not the measure-axiom level:
    \emph{$\volR$ is non-invariant under capability removal}.
    Removing an unexercised capability $c$ (a capability with
    $e_t(c) = 0$ throughout $W$) from the system leaves
    $\volR$ unchanged while $\volL$ drops by $w(c) > 0$.
    Example~\ref{ex:volR_removal_invariance} witnesses this.
    The property is about the mapping $\mathcal{P} \mapsto
    \volR^{\mathcal{P}}$ from capability systems to the measure,
    not about $\volR$ as a set function on any one
    $P^{\mathrm{ex}}_W$; it is what breaks the self-balancing
    transfer (Corollary~\ref{cor:not_self_balancing}).
  \item[(b)] \textbf{$\volRlower$: event-set non-additivity
    (not an M6 violation).}
    M6 is a measure axiom on poset-independent capability
    subsets; $\volRlower$ as a set function inherits M6 by the
    same restriction argument as M4.  A separate property
    holds at the \emph{event-set} level: merging two groups of
    sacrifice events changes per-capability attributions under
    max-attribution, and the joint maximum over all events can
    be less than the sum of within-group maxima.  This is not
    an M6 counterexample (M6 quantifies over capability subsets,
    not event sets); it is a non-superadditivity property of
    the lower-bound construction over event-set unions.
    Example~\ref{ex:m6_fail_volRlower} illustrates this property
    in a constructed scenario outside Part~B's strict admissibility
    regime, retained for operational clarity about what
    max-attribution does on merged event groups.
  \end{enumerate}
\end{description}
\end{theorem}

\begin{example}[$\volR$ is invariant under removal of
unexercised capabilities]
\label{ex:volR_removal_invariance}
Consider a single-agent system with two poset-independent
capabilities:
\begin{itemize}
\item $c_A = $ ``drive to gas station $\alpha$'', $w(c_A) = 1$.
\item $c_B = $ ``drive to gas station $\beta$'' (different
  location), $w(c_B) = 1$.
\end{itemize}
The system poset is $P = \{c_A, c_B\}$ with $\volL(P) = 2$.
$c_A$ and $c_B$ are poset-independent: different
fuel-access pathways with no shared participants.

Under the exercise indicator of
Appendix~\ref{app:exercise_indicator}, the agent uses only $c_A$
(gas station $\alpha$ is closer), so
$P^{\mathrm{ex}}_W = \{c_A\}$ and
$\volR(P) = \volL(P^{\mathrm{ex}}_W) = 1$.  (Note: $c_B$ is
\emph{not} load-bearing at any time in $W$---removing it leaves
all realized outputs unchanged.)

Now remove $c_B$: the system becomes $P' = \{c_A\}$ with
$\volL(P') = 1$ (so $\volL$ contracts by $w(c_B) = 1$).  The
agent's behavior is unchanged, so $P^{\mathrm{ex}}_W = \{c_A\}$
and $\volR(P') = 1$, unchanged.

The measure $\volR$ is therefore \emph{invariant} under the
removal of the unexercised $c_B$, while $\volL$ contracts by
$w(c_B)$.  A $\volR$-maximizing actor would face no penalty for
shedding unexercised capabilities; a $\volL$-maximizing actor
would pay $w(c_B)$.

This is not a counterexample to M6 as a measure axiom: for the
fixed exercised sub-poset $P^{\mathrm{ex}}_W = \{c_A\}$, $\volR$
trivially satisfies M6 (and all other $\volL$-inherited axioms
by restriction).  The failure is at the level of the
\emph{mapping} $\mathcal{P} \mapsto \volR^{\mathcal{P}}$: this
mapping is not \emph{strictly} monotone under capability-set
expansion ($P \supset P'$ but $\volR(P) = \volR(P')$).  Weak
monotonicity ($P \supseteq P' \implies \volR(P) \geq \volR(P')$)
holds; what fails is strict monotonicity.  Equivalently, the
mapping is not strictly contracted by capability removal when
the removed capability is unexercised.  This non-invariance is
what breaks the self-balancing transfer
(Corollary~\ref{cor:not_self_balancing}).
\end{example}

\begin{example}[$\volRlower$ non-additivity over event-set
union]
\label{ex:m6_fail_volRlower}
The example below illustrates that $\volRlower$ aggregated
across two event groups is not in general the sum of the
within-group aggregates.  This is a property of the lower-bound
construction over \emph{event sets}, not a counterexample to
M6 on capability subsets: M6 quantifies over poset-independent
capability subsets, while max-attribution operates over events
covering each capability.  We retain the example for
operational clarity (max-attribution does not behave
super-additively when event sets are merged) and flag the
distinction below.

The example is \emph{constructed outside Part~B's strict
admissibility regime}: under Part~B's within-bundle
poset-independence plus S5 plus the partition of unity
$\sum_c \alpha_{n,c} = 1$, the share-collapse argument
(Remark~\ref{rem:s5_heavy}) forces $\alpha_{n,c}$ to equal the
true share $\volR(c)/\volR(Y_n)$, leaving no room for two
events on the same bundle to disagree on $\alpha$.  We relax
those constraints below to expose the construction's
non-additivity property; the example illustrates the lower
bound's behaviour, not a Part-B-admissible scenario.  Read it
as ``if a downstream variant of the framework permitted
event-specific $\alpha$ on the same bundle (e.g., by dropping
the partition of unity), aggregating $\volRlower$ across event
groups would still fail super-additivity.''

Consider two sacrifice events with the same bundle
$Y = \{a, b\}$ and (counterfactually) different decomposition
coefficients:
\begin{itemize}
\item Event $E_1$: $\Delta\volL(X_1) = 1$, $\alpha_{1,a} = 0.5$,
  $\alpha_{1,b} = 0.5$.
\item Event $E_2$: $\Delta\volL(X_2) = 1$, $\alpha_{2,a} = 0.3$,
  $\alpha_{2,b} = 0.7$.
\end{itemize}
Let group $A = \{E_1\}$ and group $B = \{E_2\}$.  Under
Theorem~\ref{thm:aggregate_bound} Part~B:
\begin{align*}
  \volRlower(a \mid A) &= 0.5, &
  \volRlower(b \mid A) &= 0.5, \\
  \volRlower(a \mid B) &= 0.3, &
  \volRlower(b \mid B) &= 0.7.
\end{align*}
Summing per-capability:
$\volRlower(A) = 1.0$, $\volRlower(B) = 1.0$.

In the merged group $A \cup B$, max-attribution over both events:
$\volRlower(a \mid A \cup B) = \max(0.5, 0.3) = 0.5$ and
$\volRlower(b \mid A \cup B) = \max(0.5, 0.7) = 0.7$.
Total: $\volRlower(A \cup B) = 1.2$.

If we (informally) compare a within-group sum
$\volRlower(A) + \volRlower(B) = 2.0$ to the merged-event
aggregate $\volRlower(A \cup B) = 1.2$, the merged aggregate
is strictly smaller.  This is \emph{not} an M6 counterexample
on capability subsets: $A$ and $B$ here are event sets, not
poset-independent capability subsets, and M6 makes claims
about the latter, not the former.  What the example shows is
that \emph{max-attribution does not commute with event-set
union}: merging two event groups whose attributions disagree
on the same capability produces a per-capability max bounded
above by the sum of within-group maxes.  When restricted to
the M6 setting (poset-independent capability subsets, with
each capability covered by a single attribution event), this
non-additivity does not arise; M6(b) below therefore fails for
$\volRlower$ specifically through the event-set-union
subadditivity exhibited here, not through any axiom-level
failure on capability subsets.

This is a property of the \emph{lower-bound construction over
event sets}, not of the underlying $\volR$: the true $\volR$
aggregate of the merged events need not be subadditive, but
our observational lower bound is.  The non-self-balancing
corollary below therefore relies on the capability-removal
non-invariance of M6(a), not on this lower-bound
event-set subadditivity.
\end{example}

\begin{corollary}[$\volR$ Is Not Self-Balancing]
\label{cor:not_self_balancing}
The self-balancing property does \emph{not} transfer from
$\volL$ to $\volR$.  Paper~2's self-balancing
theorem~\citep[Theorem~1]{lasser2026poset} relies on the fact
that \emph{any} capability removal contracts $\volL$, making
$\volL$-maximization incompatible with unilateral capability
reduction.  For $\volR$, removing an unexercised capability
leaves the measure unchanged
(Example~\ref{ex:volR_removal_invariance}), so a
$\volR$-maximizer can shed unused capabilities without penalty.
The protective coupling between measure-maximization and
capability preservation that $\volL$'s self-balancing encodes
therefore does not apply to $\volR$.

The corollary does \emph{not} follow from $\volR$ failing M6 as
a measure axiom---it does not
(Theorem~\ref{thm:axiom_inheritance}, M6(a))---nor from
$\volRlower$ failing M6 alone (a lower bound failing measure
axioms is unremarkable).  The corollary rests on the
system-level non-invariance of M6(a): $\volR$'s value depends
on which capabilities are exercised, and unexercised
capabilities can be removed without affecting that value.
\end{corollary}

This is the structural reason to use $\volR$ as a diagnostic
rather than an objective.  An actor maximizing $\volR$ would
lack the automatic diversity-preservation that makes $\volL$
safe: exercise is rewarded, but dormant optionality is not
protected.

\begin{remark}[When M6 does hold for $\volRlower$]
M6 holds for $\volRlower$ when the merged groups have sacrifice
events in poset-independent bundle categories
(Definition~\ref{def:category_partition})---each unbenchmarked
capability appears in sacrifice events from only one group, and
cross-group capabilities are poset-independent.  Under this
non-redundancy condition, $\volRlower$ inherits M1--M6.  Even
so, the self-balancing property transfers to $\volRlower$ only as
a \emph{lower bound} on self-balancing---$\volRlower$ being
self-balancing does not imply $\volR$ itself is self-balancing.
\end{remark}
